\documentclass[a4paper,11pt]{article}
\usepackage[margin=1in]{geometry}
\usepackage{hyperref}
\usepackage{parskip}
\urlstyle{rm}
\pagestyle{empty}

\begin{document}

\textbf{\Large Shubham Pandey}\\
\href{mailto:mr.pshubham@gmail.com}{mr.pshubham@gmail.com} \quad +91-8174081896\\
\href{https://www.linkedin.com/in/shubham-pandey-a74716199/}{linkedin.com/in/shubham-pandey-a74716199} \quad
\href{https://github.com/shubham-pandey12}{github.com/shubham-pandey12}

\vspace{1em}
October 30, 2026

\vspace{1em}
Hiring Team\\
Chanakya Team\\
Sarvam AI\\
Bengaluru, India

\vspace{1em}
Dear Hiring Team,

The Frontend Engineer -- Chanakya role at Sarvam AI is a strong match for the frontend work I do --- building production-grade, component-driven interfaces for enterprise and high-traffic environments, owning full frontend delivery from architecture to deployment, and applying performance engineering rigor to constrained environments.

At Moresand, I design and own performant, production-grade frontend interfaces (React.js, TypeScript, Tailwind CSS) for a multi-tenant SaaS platform --- building reusable component libraries, shared state management (Redux), and async API patterns across 5+ enterprise client deployments. I apply performance engineering (lazy loading, rendering optimization) and own the complete frontend layer: from component architecture through CI/CD to production monitoring. At MAQ Software, I built production AI-enabled data dashboards (React.js, TypeScript) with real-time async data flows and REST API integration, owning the frontend delivery end-to-end with Claude and GitHub Copilot. My WebHealth project uses \textbf{Recharts} for data visualization dashboards and \textbf{Three.js} for 3D visual reporting --- directly matching your data visualization requirements, with Google Lighthouse integration for performance engineering. At GfG, I optimized React.js and Next.js components on a platform serving millions of monthly users, improving performance by 40\%.

Sarvam's focus on building atoms --- replicable, deployable product components --- for demanding, constrained environments is exactly the kind of ownership I find compelling. Building interfaces that work offline-first, in low-bandwidth environments, for India's most demanding institutional deployments is a frontier engineering problem. I am excited to contribute to a team building Sovereign AI infrastructure for India at this scale.

\vspace{1em}
Sincerely,\\[0.5em]
\textbf{Shubham Pandey}

\end{document}
